\chapter{Morality and grammaticality}

\section{Ideas}
\begin{enumerate}
    \item The Pirahã
    \item Alice Evans and \href{https://www.youtube.com/watch?v=6gC8Mi_sOQU}{link}
    \item A difference between moral intuitions and grammatical intuitions is that it is typically less self-serving to adjust your grammatical intuitions. (true?)
\end{enumerate}

Could speech have real-world negative consequences by virtue of its being ungrammatical?
When Eve eats of the fruit of the tree of knowledge of good and evil, she already speaks.

\section{Utilitarianism}

Utilitarianism is \dots

Bentham's musings on language, writing, and grammar offer another dimension to his character. Living in an era where legal and philosophical texts were often mired in complexity, Bentham championed clarity and precision. He was particularly critical of the legal jargon of his time, viewing it as a barrier to justice and understanding. ``Prolixity is the soul of tedium,'' he once remarked, emphasizing the need for simplicity in expression. For Bentham, language was a tool, its value derived from its utility in conveying ideas. This pragmatic view extended to his thoughts on grammaticality. While he acknowledged the importance of grammatical norms, he believed that the ultimate aim of language was clear communication. Any rule that hindered this objective was, in his view, superfluous.

\subsection{Language frameworks}

Many language assessment schemes, such as the Canadian Language Benchmarks, or the Common European Framework of Reference for Language take a pragmatic, utilitarian view of grammaticality. Language proficiency, under these schemes, is very much about what you can achieve, what you can do, what your instrumental outcomes are when communicating in the language. For example, at CLB 7 (out of 12), ``Adequate Intermediate Ability'', \begin{quote}
    The speaker can: Communicate with some confidence in many daily routine social, educational, and work situations, and present concrete and some abstract information on an expanding range of familiar topics.
\end{quote} And in so doing, the speaker demonstrates an \begin{quote} 
    increasing variety of grammatical structures, with developing control of complex structures, [although] grammar, vocabulary and pronunciation difficulties may occasionally impede communication.
\end{quote}

\subsection{Cooperation and grammar}

Jahangir Khan is anything but cooperative. A titan in the world of squash, he dominated the sport in the 1980s. His unparalleled winning streak of 555 matches remains unmatched in professional sports. Khan's style of play was relentless, often wearing down his opponents with his exceptional fitness and tactical acumen. Yet, his success was not built on cooperation. In the court's confines, he was a fierce competitor, always seeking an edge, always looking for a way to outplay his adversary. His non-cooperative stance in the game was not a sign of moral failing but a strategic choice, essential for triumph in such a competitive arena.

Jahangir Khan's life and career are replete with instances that showcase his resilience, determination, and humility. Born into the squash-playing Khan family of Pakistan, Jahangir's introduction to the sport was almost preordained. But his journey was anything but smooth. At a young age, he faced health challenges, including a couple of hernia operations. Few believed he would ever play professionally, let alone dominate the sport.

Tragedy struck when Jahangir's older brother, Torsam Khan, died suddenly during a game due to a heart attack. Torsam had been a guiding force in young Jahangir's life, mentoring him in the nuances of the game. Devastated by the loss, Jahangir considered quitting the sport, but he chose to persevere, turning his grief into motivation. He once remarked that he felt his brother's presence with him during matches, guiding and watching over him.

Off the court, Jahangir was known for his humility. Despite his unmatched success, he was always approachable and took time to engage with fans and budding players. There are numerous accounts of him offering tips to young players and always being willing to share his knowledge of the game. His sportsmanship was also legendary. 

Khan was clearly a deeply moral and ethical person, but that didn't prevent him from intense competition on the court. Many spheres of human interaction are built on the bedrock of cooperation. Consider the intricate dance of a symphony orchestra, where each musician must not only master their instrument but also listen intently to their fellow players, synchronizing their efforts to produce harmonious sound. Here, cooperation isn't just beneficial; it's essential.

Yet, the principle of cooperation isn't universally applicable. Just as Khan found success in non-cooperation on the squash court, there are contexts in life where competition is crucial: sports, business, arguing legal cases, science. This is not to say that all interactions in those fields are competitive -- a lawyer has to cooperate with their client -- but there are particular interactions where cooperation is inimical to the proper functioning of the institution.

going against the grain, challenging the status quo, or simply choosing a different path can lead to innovation and progress. Think of the mavericks in the world of art or science who defied conventions and, in doing so, reshaped entire fields.

Language, too, offers a rich tapestry of cooperative and non-cooperative dynamics. At its core, communication is a cooperative endeavor. Speakers and listeners engage in a mutual dance, with each party making assumptions, drawing inferences, and adjusting their behavior based on feedback. The Gricean maxims, which emphasize the importance of being clear, relevant, truthful, and concise, underscore this cooperative nature of communication.

However, just as in squash or music, there are times when breaking the rules or flouting conventions in language can be powerful. Think of the poets, playwrights, and novelists who play with words, bending and twisting them to create new meanings or evoke specific emotions. Their deviations from linguistic norms aren't errors but deliberate choices, much like Khan's tactical decisions on the court.

This brings us to the concept of grammaticality. In certain circles, especially academic or formal settings, adhering to the norms of Standard English is expected. But these norms are not set in stone. They evolve over time, shaped by societal changes, technological advancements, and cultural shifts. What's deemed ``ungrammatical'' today might be the standard of tomorrow.

Moreover, the idea of a ``standard'' itself is fraught with complexities. Who decides what's standard? Why are certain dialects or forms privileged over others? These questions hint at the power dynamics at play in language, where issues of prestige, class, and identity intersect with grammar.

In everyday interactions, people often switch between different linguistic registers, choosing the one that's most appropriate for the situation. A casual chat with a friend might see the use of slang, colloquialisms, and non-standard forms, while a job interview would demand more formal language. This linguistic flexibility is a testament to our ability to navigate different social contexts, recognizing when to cooperate with linguistic norms and when to deviate.

Returning to Khan, his non-cooperative strategy in squash was a response to the specific demands of his sport. Similarly, our linguistic choices are shaped by the demands of our communicative context. Recognizing when to adhere to linguistic norms and when to challenge them is a skill, honed over time and through experience.

In conclusion, while cooperation is a guiding principle in many areas of life, it's not an absolute. Just as Khan's dominance in squash was built on a foundation of strategic non-cooperation, there are times in language, and life more broadly, when going against the norm can be not just beneficial, but transformative.

The key lies in discernment, in understanding the nuances of each situation, and in making informed choices. Whether on the squash court, in a concert hall, or in a conversation, recognizing the interplay of cooperation and non-cooperation is crucial.

\section{Laws and morality}

Moral philosophers will sometimes point out a distinction between laws and moral

\subsection{Double modals in English}

Linguists, like bird watchers eagerly seeking rare bird species, are fascinated by unusual language patterns. One such rare pattern in English is the double-modal construction, exemplified by phrases like (\ref{ex:would-might}).

\ea{\textit{The flavouring within soda can come from a wide variety of places, but one you \uline{would might} not expect would be from tree bark.}\footnote{\href{https://youtu.be/_bAnbHsDWa0?feature=shared&t=157}{YouTube}}}\label{ex:would-might}
\z 

I can't ever recall hearing a single instance of a double modal in my daily life, but Steven Coats, of University of Oulu in Finland, got tired of waiting and went looking. He started by developing the 1.29-billion-word \textit{Corpus of North American Spoken English} \citep{Coats2023} from YouTube videos. Then, he searched this massive corpus for double modals like \textit{would might}. Next, he had to check each one and determine whether it was a real double modal or just some kind of transcription error, self correction, change of speaker, subordinate clause (\textit{those of us }[\textit{who can}] [\textit{must make ...}]) or any number of other false positives. After all that, just 54 true example of \textit{would might} remained \citep{Coats2022}. That's about one for each 24 million words. For comparison, \textit{interstitial} occurs about once per 2 million words, while you'd have to go to a word like \textit{exigence} to find one that is roughly as rare. Coats tells me he finds just three in my home province of Ontario.

For most English speakers, double modals are ungrammatical. Not only are they ungrammatical, but they are so rare in most communities that most of us are entirely unaware that they even exist. We've probably never heard them deprecated, unlike constructions such as \textit{him and me tried it}, \textit{my and Geoff's book}, \textit{I ain't gonna try that}, \textit{I done it already}, and a host of others.

There are a good grammatical reason for this. In Standard English, the modal auxiliary verbs (\textit{may}, \textit{might}, \textit{can}, \textit{could}, \textit{shall}, \textit{should}, \textit{will}, \textit{would}, and \textit{must}, along with their negative forms, \textit{can't}, \textit{won't}, etc.) are highly defective verbs. To call a verb defective is not to render some qualitative or moral judgment, but rather to indicate that it lacks one or more of the regular verb forms. A normal verb like \textit{jump} has six forms: a plain form (\textit{jump}), a plain present-tense form (also \textit{jump}), a third-person singular present-tense form (\textit{jumps}), a past-tense form (\textit{jumped}), a past participial form (also \textit{jumped}), and a present participial form (\textit{jumping}).

Contrast that with a modal auxiliary verb like \textit{will}, with only four: present-tense \textit{will}, past-tense \textit{would}, negative present-tense \textit{won't}, and negative past-tense \textit{wouldn't}. The modal auxiliaries lack a plain form, the form used in infinitival clauses (\textit{It makes you \uline{jump}.}) You can't say *\textit{It makes you will jump} because \textit{will} is the present-tense, and \textit{make} requires the plain form. If a present-tense form worked, \textit{jumps} would be fine in *\textit{it makes him jumps}, and it's not. Other places the require the plain form are after \textit{to} (\textit{I want to }(*\textit{can})\textit{ jump}), after \textit{let's} (\textit{Let's }(*\textit{can})\textit{ jump}), and -- here's the crucial part -- after a modal auxiliary verb (\textit{I should }(*\textit{can})\textit{ jump}). In Standard English, you don't find \textit{you \uline{would might} not expect it} because \textit{would} expects a non-tensed verb, and \textit{might} is past-tense.

\section{Ethical Perspectives on Grammaticality}

Grammaticality judgments, often presented as objective facts about language, frequently reflect underlying ethical stances. These judgments intersect with social values, power structures, and cultural identities in complex ways.

Derrick Parfit's late-career effort to unify disparate moral theories offers a useful lens through which to view these varied approaches to grammaticality. Parfit's ``climbing the mountain'' metaphor suggests that seemingly conflicting ethical frameworks might converge at a higher level of understanding. Similarly, our varied perspectives on grammaticality might be seen as different paths up the same mountain.

Consider how grammatical "errors" are often treated as moral failings, or how grammatical norms typically align with the speech of socially dominant groups. These phenomena suggest that our notions of grammaticality extend beyond mere linguistic description into the realm of ethics and social justice.

The following sections explore seven ethical frameworks and their implications for understanding grammaticality. Each offers a distinct perspective on what makes language use "right" or "wrong." Yet, following Parfit's insight, we might find that these perspectives are not as irreconcilable as they first appear.

This exploration has real-world implications. Our conceptualization of grammaticality shapes language policy, educational practices, and social attitudes. By seeking a unified understanding of grammaticality's ethical dimensions, we might develop more nuanced and effective approaches to language variation and change in our diverse linguistic landscape.

\section{Ethical Perspectives on Grammaticality}

Grammaticality judgments are often presented as objective facts about language. Yet these judgments frequently reflect deeper, often unexamined, ethical stances. The connection between grammaticality and morality isn't immediately obvious, but it becomes clearer when we consider how language norms intersect with social values, power structures, and cultural identities.

Consider how grammatical ``errors'' are often treated not just as mistakes, but as moral failings. A misplaced apostrophe or a split infinitive might elicit reactions of disgust or contempt far out of proportion to any communicative impairment they cause. This suggests that our notions of grammaticality go beyond mere linguistic description.

Grammatical norms often align with social power structures. Standard language varieties typically correspond to the speech of socially dominant groups. This raises questions about the ethical implications of enforcing these norms. Are we simply maintaining useful standards for clear communication, or are we perpetuating social inequalities?

These issues become particularly salient in contexts like education, where judgments about language can have profound effects on students' opportunities and self-perception. They also play out in professional settings, where linguistic choices can influence hiring decisions and career advancement.

By examining grammaticality through various ethical lenses, we can gain insight into these complex issues. This approach allows us to question assumptions about what constitutes ``correct'' language use and to consider the broader implications of how we define and enforce grammatical norms.

The following sections explore seven ethical frameworks and their implications for our understanding of grammaticality. Each offers a distinct perspective on what makes language use ``right'' or ``wrong,'' and each has different consequences for how we approach language in practice.

This exploration is not merely academic. How we conceptualize grammaticality shapes language policy, educational practices, and social attitudes. By understanding the ethical dimensions of grammaticality, we can make more informed decisions about how to approach language variation and change in our increasingly diverse linguistic landscape.

\subsection{Pragmatism: Effectiveness Over Prescription}

Pragmatism in grammaticality prioritizes effective communication over strict adherence to prescriptive rules. This approach aligns closely with the utilitarian view of language proficiency exemplified by frameworks such as the Canadian Language Benchmarks (CLB) or the Common European Framework of Reference for Language (CEFR).

Consider the CLB description for level 7, ``Adequate Intermediate Ability'':

\begin{quote}
The speaker can: Communicate with some confidence in many daily routine social, educational, and work situations, and present concrete and some abstract information on an expanding range of familiar topics.
\end{quote}

Notably, this description focuses on communicative outcomes rather than perfect grammatical form. It acknowledges that at this level, a speaker demonstrates an ``increasing variety of grammatical structures, with developing control of complex structures," while recognizing that ``grammar, vocabulary and pronunciation difficulties may occasionally impede communication."

This pragmatic stance echoes Jeremy Bentham's views on language. Bentham, living in an era of often impenetrable legal and philosophical texts, championed clarity and precision. For Bentham, language was a tool, its value derived from its utility in conveying ideas. While he acknowledged the importance of grammatical norms, he believed that the ultimate aim of language was clear communication. Any rule that hindered this objective was, in his view, superfluous.

The pragmatic approach to grammaticality raises interesting questions about the nature of linguistic cooperation. Paul Grice's cooperative principle and its associated maxims (quality, quantity, relevance, and manner) suggest that effective communication relies on a kind of tacit agreement between speakers and listeners. However, as we've seen with examples like Jahangir Khan in squash, non-cooperation can sometimes be strategically advantageous.

In language, this non-cooperative stance might manifest as deliberate deviations from grammatical norms for rhetorical effect. Consider the following examples:

\begin{enumerate}
    \item Poetic license: "Anyone lived in a pretty how town" (E.E. Cummings)
    \item Advertising slogans: "Think different" (Apple)
    \item Political rhetoric: "We shall fight on the beaches, we shall fight on the landing grounds, we shall fight in the fields and in the streets, we shall fight in the hills" (Winston Churchill)
\end{enumerate}

Each of these examples breaks conventional grammatical rules, yet each is highly effective in its context. From a pragmatic perspective, they could be considered grammatical because they successfully achieve their communicative goals.

\bigskip


The pragmatic approach, though, is not without challenges. One significant issue is the potential conflict between immediate communicative effectiveness and broader societal needs for standardization. While "u ok?" might be a perfectly effective text message, accepting such forms in all contexts could lead to communication difficulties across different groups or in more formal situations.

Moreover, judgments of communicative effectiveness can be subjective and culturally bound. What constitutes clear communication in one context might be opaque or inappropriate in another. This subjectivity can make it difficult to establish consistent standards, particularly in educational or professional settings.

Consider the case of double modals in English, such as "might could" or "used to could." These constructions, while rare in most varieties of English, are grammatical in some Southern United States dialects. A strictly pragmatic approach might accept these forms as grammatical wherever they're understood, but this could lead to difficulties in contexts where speakers are unfamiliar with the construction.

The pragmatic approach to grammaticality also intersects with questions of linguistic identity and power. Non-standard forms often serve as markers of group identity or solidarity. Judging these forms solely on their communicative effectiveness to outsiders might miss their important social functions within the community.

In conclusion, while the pragmatic approach offers valuable insights into the nature of grammaticality, it cannot be applied in isolation. A comprehensive understanding of grammaticality must balance communicative effectiveness with considerations of standardization, cultural context, and the social functions of language variation. As with many aspects of language, the key lies in recognizing the complex interplay between these factors rather than seeking a single, universal criterion for grammaticality.

\subsection{Virtue Ethics: Cultivating Linguistic Character}

Virtue ethics shifts our focus from rules or outcomes to the character of the language user. It's less about whether a particular utterance is right or wrong and more about developing linguistic virtues like clarity, precision, and eloquence.

This approach aligns with how we often judge writing. We don't just look at whether the grammar is correct, but at how well the author uses language to convey their ideas. A virtuosic writer isn't just following rules; they're wielding language with skill and insight.

But what does this mean for everyday language use? Let's consider an example:

\begin{quote}
A: ``The data don't support your conclusion.''\\
B: ``The data doesn't support your conclusion.''
\end{quote}

Prescriptively, A is correct (``data'' is plural). But B is increasingly common, even among educated speakers. A virtue ethics approach might ask: Which speaker is exhibiting better linguistic character? Is A showing precision and respect for traditional forms? Or is B demonstrating adaptability to changing usage?

There's no clear-cut answer. It depends on context, audience, and the speaker's overall command of language. A scientist writing a paper might lean towards A, showing mastery of field-specific conventions. A politician addressing a general audience might prefer B, displaying an ability to communicate clearly without sounding pedantic.

This perspective encourages us to view grammaticality not as a binary state but as a skill to be honed. It's about developing judgment about when to adhere strictly to rules and when flexibility serves communication better.

Consider how we judge a chef. We don't just check if they followed a recipe correctly. We look at how they combine ingredients, adapt to available materials, and create dishes that suit their diners. Similarly, a virtue ethics approach to grammaticality looks at how speakers combine linguistic elements, adapt to communicative contexts, and craft utterances that serve their audience.

This view of grammaticality as a cultivated skill rather than a set of inflexible rules echoes how linguists understand language acquisition. Children don't learn their native language by memorizing rules. They develop an intuitive grasp of the language's patterns through exposure and use.

The virtue ethics approach also offers an interesting perspective on language change. If we see grammaticality as a matter of cultivated skill rather than rigid rules, it becomes easier to understand how and why languages evolve. Changes aren't ``errors'' that have become widespread, but shifts in what the language community considers skillful usage.

Take the ongoing change in how we use ``literally.'' Prescriptivists bemoan its use as an intensifier (``I literally died laughing''), but this usage is becoming increasingly accepted. A virtue ethics approach might say that skillful use of language now includes knowing when this intensifier use of ``literally'' is appropriate and when it isn't.

This perspective on grammaticality isn't without challenges. It's subjective and context-dependent, which can make it difficult to teach or assess. It also risks reinforcing linguistic prejudices if we're not careful about whose ``virtues'' we're promoting.

But despite these challenges, a virtue ethics approach to grammaticality offers a nuanced way of thinking about language use. It encourages us to move beyond simple right/wrong judgments and consider the broader skills involved in effective communication. 

\subsection{Consequentialism: Judging Language by Its Outcomes}

Consequentialism in grammaticality focuses on the results of language use rather than adherence to rules. It's about whether an utterance achieves its intended effect, not whether it follows prescribed norms.

This view aligns with how we often judge communication in real-world contexts. In business, for example, we care more about whether a message persuades a client than if it splits an infinitive. In casual conversation, we're more concerned with understanding and connection than perfect subject-verb agreement.

Let's look at an example:

"Me and him went to the store."

Prescriptively, this is incorrect. It should be "He and I went to the store." But a consequentialist approach asks: Did the message get across? Did any misunderstanding occur? If not, then from this perspective, the sentence worked. It was grammatical enough.

This doesn't mean "anything goes." Consequentialism in grammar isn't about abandoning all standards. It's about judging those standards by their effects. Some prescriptive rules genuinely aid clarity and prevent misunderstanding. Others don't. Consequentialism gives us a framework for distinguishing between them.

Consider the "rule" against ending sentences with prepositions. Does following this rule actually improve communication? Often, it doesn't. It can lead to awkward constructions like "This is the sort of arrant pedantry up with which I will not put." A consequentialist view might say this rule isn't worth following most of the time because its outcomes aren't positive.

On the other hand, the rule about using "fewer" for countable items and "less" for uncountable ones can prevent ambiguity. "Fewer calories" and "less calorie intake" convey different meanings. Here, the rule serves a purpose. It leads to clearer communication.

The consequentialist approach to grammaticality also offers insight into language change. New forms that effectively serve communicative needs tend to stick around and spread. Those that don't, don't. This is why "u" for "you" in text messages has caught on, but most of us aren't dropping our definite articles despite their frequent omission in headlines.

This view of grammaticality isn't without problems. It can be hard to measure the consequences of language use objectively. It might also undervalue the social and cultural roles of language that go beyond mere information transfer. But it offers a practical framework for thinking about grammaticality, one that aligns well with how we often judge language in day-to-day life.

\subsection{Deontology: Grammar as a Moral Duty}

Deontology in grammaticality posits that following grammatical rules is a duty, regardless of outcomes. This view aligns closely with prescriptivism in linguistics.

For a deontologist, the sentence "Me and him went to the store" is wrong not because it fails to communicate or because it leads to bad consequences, but simply because it breaks the rules of Standard English grammar. The correct form, "He and I went to the store," is right because it follows the rules.

This approach has some appeal. It offers clear, consistent standards. It aligns with how grammar is often taught in schools. And it reflects a common intuition that there's a "right" and "wrong" way to use language.

But deontology in grammar faces significant challenges:

\begin{enumerate}
    \item Which rules? English has many varieties, each with its own norms. Even within Standard English, experts often disagree on specific rules.
    
    \item Rule conflicts: What happens when rules contradict each other? The injunction against splitting infinitives often clashes with the desire for clear, concise expression.
    
    \item Language change: Deontology struggles to account for how language evolves. Today's error often becomes tomorrow's standard usage.
    
    \item Linguistic diversity: A strict rule-based approach can delegitimize non-standard dialects, even when they're systematic and effective for their speakers.
\end{enumerate}

Consider the double modal construction we discussed earlier:

\ea \textit{The flavouring within soda can come from a wide variety of places, but one you \uline{would might} not expect would be from tree bark.}
\z

For most English speakers, this is ungrammatical. The deontological view would say it's wrong, full stop. But this ignores the fact that double modals are grammatical in some dialects. It also misses the potential communicative value of the construction, which can convey nuances of meaning not easily expressed with single modals.

Deontology in grammar often manifests as ``grammar purism,'' the belief that there's one correct way to use language and all deviations are errors. This view can lead to linguistic prescriptivism and even language discrimination.

Yet, the deontological approach isn't without value. Standardized rules can facilitate clear communication, especially in formal or professional contexts. They can provide a shared framework for language learners. And for many people, following grammatical rules is part of showing respect for the language and its traditions.

The key is recognizing the limits of this approach. Rules can guide us, but they're not absolute. They're tools for communication, not moral imperatives. A nuanced view of grammaticality needs to balance respect for linguistic norms with an understanding of language's inherent flexibility and diversity.

\subsection{Elitist Power: Grammar as Social Currency}

The elitist power view of grammaticality treats ``correct'' language use as a marker of social status and authority. Under this framework, adherence to prescriptive grammar rules isn't just about clear communication -- it's a way to signal membership in educated or elite social groups.

But it's not just elites who use language to discriminate. People across the social spectrum create and use non-standard forms, particularly slang, to set themselves apart. This linguistic differentiation serves to strengthen in-group bonds and exclude outsiders.

Consider the reaction to a sentence like:

\ea
\textit{Me and John was talking about that yesterday.''}
\z

An elitist approach might not just label this as ungrammatical, but use it to make judgments about the speaker's education, class, or fitness for certain social roles. The ``correct'' version -- \textit{John and I were talking about that yesterday} -- becomes a shibboleth, a way to identify who belongs and who doesn't.

This perspective often manifests in the enforcement of grammatical ``rules'' that have little to do with clarity or effective communication. The prohibition against split infinitives, for instance, originated as a way for educated elites to signal their knowledge of Latin grammar. It persists not because it aids understanding, but because it serves as a social marker.

The elitist view of grammaticality poses several problems:

\begin{enumerate}
    \item It often privileges the dialect of dominant social groups, treating linguistic differences as deficiencies.
    \item It can be used to justify discrimination in education, employment, and other areas.
    \item It ignores the linguistic validity and communicative effectiveness of non-standard forms.
    \item It tends to resist language change, treating evolving usage as ``decay'' rather than natural evolution.
\end{enumerate}

Yet this approach persists, often implicitly, in many institutions. School systems that penalize students for using non-standard dialects, workplaces that judge employees' competence based on their mastery of prescriptive grammar, and social circles that use language as a gatekeeping mechanism all reflect aspects of this elitist view.

It's worth noting that linguistic identity and social identity are deeply interconnected. Speakers of prestige dialects are unlikely to broadly adopt elements of non-pre\-stige dialects, though certain ``cool'' forms may be ``appropriated.'' This interconnection makes changing linguistic hierarchies particularly challenging.

\subsection{Social Justice by Any Means Necessary: Grammar as Liberation}

This radical approach to grammaticality views traditional grammatical norms as tools of oppression and seeks to dismantle them. It's less about defining what's grammatical and more about questioning why certain forms are privileged over others.

Consider these lines from ``Fight the Power'':

\ea
\textit{Gotta give us what we need\\
Our freedom of speech is freedom or death\\
We got to fight the powers that be}
\z

Standard grammar labels this ungrammatical. But a social justice approach might argue that this form is valid and that labelling it ``incorrect'' perpetuates linguistic discrimination. It recognizes that these non-standard forms are deliberately chosen for their expressive power and as markers of group identity.

This perspective often manifests in:

\begin{itemize}
    \item Advocating for acceptance of non-standard dialects in a broader range of contexts, including formal contexts
    \item Criticizing the notion of a single ``standard'' language
    \item Promoting ``code-meshing'' (blending dialects) over ``code-switching''
    \item Questioning the cultural biases in grammar textbooks and standardized tests
\end{itemize}

The social justice view of grammaticality faces several challenges:

\begin{enumerate}
    \item It can sometimes overlook the practical value of a shared linguistic standard
    \item It might inadvertently reinforce linguistic isolation of marginalized groups
    \item It struggles to define boundaries -- if all forms are equally valid, how do we maintain mutual intelligibility?
    \item It can be seen as threatening by those who've invested heavily in mastering standard forms
\end{enumerate}

Interestingly, the elitist power view and the social justice view can be seen as mirror images of each other, with the groups reversed. While the elitist view seeks to maintain linguistic hierarchies that benefit dominant groups, the social justice view aims to dismantle these hierarchies and elevate non-standard forms. Both approaches use language as a tool for group identity and social positioning, albeit with very different goals and beneficiaries.

This mirroring extends to the challenges each view faces. Just as the elitist view can reinforce social inequalities, the social justice approach risks isolating speakers of non-standard dialects if taken to extremes. Both views also grapple with the tension between group identity and broader communication needs, albeit from opposite directions.

Moreover, both approaches often align with identity politics, which can take a zero-sum view of power and prestige. While economic resources can be expanded, it's less clear that linguistic prestige can be similarly enlarged. This raises questions about the feasibility of truly equalizing linguistic status across all varieties.

There's also a tension between the desire for linguistic liberation and the practical reality that certain forms of linguistic capital are more valued in many social and professional contexts. While advocating for broader acceptance of diverse language forms, this approach must grapple with the fact that mastery of standard forms often remains a pathway to social mobility.

Yet the social justice approach offers valuable insights. It highlights how grammatical ``rules'' often reflect and reinforce social power structures. It encourages us to examine our linguistic biases and consider the rich diversity of language use across communities.

Consider the ``singular \textit{they}.'' Long criticized as ungrammatical, it's gaining acceptance partly due to social justice arguments about gender inclusivity. This shows how this approach can drive language change.

Ultimately, both the elitist power view and the social justice view of grammaticality push us to think critically about language norms, albeit in very different ways. They serve as counterbalances to each other, encouraging a more nuanced understanding of what counts as ``grammatical'' and how language intersects with social power structures. However, both must also reckon with the complex realities of linguistic identity and the practical implications of linguistic diversity in a world that often demands standardization.

\subsection{Conclusion: A Unified Theory of Grammaticality}

Our exploration of ethical perspectives on grammaticality has revealed approaches that often seem at odds. Yet, Parfit's work on unifying ethical theories offers a model for reconciling these varied views on language norms. Just as Parfit argued that consequentialism, Kantian deontology, and contractualism are ``climbing the same mountain," we can view our grammatical frameworks as different routes to a unified understanding of language norms.

This unified theory of grammaticality would blend elements from each perspective:

\paragraph{Rule Convergence} Different grammatical approaches often yield similar judgments in practice, much as Parfit proposed that seemingly different moral rules often recommend the same actions. For instance, both prescriptivists and pragmatists might prefer standard forms in formal writing, albeit for different reasons. Our unified theory would recognize these convergences, focusing on the common ground rather than surface disagreements.

\paragraph{Underlying Principles} Across grammatical perspectives, we can identify shared concerns: clarity, social cohesion, identity expression, and adaptability to change. Our unified theory would articulate these common principles, showing how they manifest in different approaches to grammaticality.

\paragraph{Triple Theory} Parfit's "Triple Theory" combined Kantian ethics, consequentialism, and contractualism. A unified theory of grammaticality would similarly blend rule-based, outcome-focused, and social contract elements of language norms by:

\begin{itemize}
    \item Acknowledging the role of conventional rules while recognizing their historical contingency
    \item Evaluating grammatical choices based on their communicative outcomes in specific contexts
    \item Treating language norms as implicit social agreements that can be renegotiated as society changes
\end{itemize}

For example, consider the sentence "Me and him went to the store." Our unified approach would:

\begin{itemize}
    \item Recognize the prescriptive rule preferring "He and I went to the store"
    \item Assess whether the non-standard form impedes understanding in its context
    \item Consider how acceptable this form is within the relevant speech community
\end{itemize}

\paragraph{Wide Value-Based Objective Theory} Parfit argued for moral realism based on irreducibly normative truths. Similarly, our unified theory would posit that there are objective truths about language use, but that these are complex, context-dependent, and not reducible to simple prescriptive rules. It would recognize that grammaticality judgments involve both universal linguistic principles and culturally specific norms.

This unified approach views grammaticality not as a binary state, but as a multidimensional space where linguistic, social, and ethical factors intersect. It acknowledges the validity of different perspectives while seeking common ground.

In practice, this theory would lead to more nuanced language policies and pedagogies. Instead of teaching absolute rules or embracing complete relativism, it would foster critical language awareness. Language users would learn to navigate the complex landscape of grammaticality, making informed choices based on context, audience, and communicative intent.

For instance, in language education, this approach might:

\begin{itemize}
    \item Teach standard forms while explaining their social and historical context
    \item Explore how grammatical choices affect communication in various settings
    \item Discuss how language norms reflect and shape social relationships
    \item Encourage students to analyze and question grammatical judgments
\end{itemize}

In language policy, it could lead to:

\begin{itemize}
    \item More flexible standards that account for linguistic diversity
    \item Policies that balance the need for a common standard with respect for non-standard varieties
    \item Recognition of the legitimacy of multiple language varieties in different contexts
\end{itemize}

Challenges remain, of course. Balancing standardization with linguistic diversity, reconciling descriptive and prescriptive impulses, and addressing ingrained language attitudes are ongoing tasks. But by recognizing the common ground among different approaches to grammaticality, we can work towards a more comprehensive, ethically informed understanding of language norms.

This unified theory doesn't eliminate the tensions between different perspectives on grammaticality. Rather, it provides a framework for understanding these tensions and navigating them productively. Like Parfit's approach to ethics, it seeks not to declare one view correct and others wrong, but to show how seemingly conflicting views can be seen as different aspects of a more complex, multifaceted truth about language and its role in human society.

\section{Synthesis: Ethical Frameworks and Linguistic Realities}

Having examined grammaticality through the lenses of consequentialism, deontology, and contractualism, we now turn to synthesizing these perspectives and grounding them in linguistic realities. This synthesis reveals both the strengths and limitations of applying ethical frameworks to language, while also highlighting crucial factors in grammaticality that these frameworks might overlook.

\subsection{Comparing Ethical Approaches}

